\section{Does a nonlinear bispectrum give a more realistic prediction?}
\label{sec:nonlinear}

This is the question the programme started from. The comparison run here is
tree level against BiHalofit, on the corrected grid, as a function of the
multipole cutoff, so that the model change and the convergence question are
not confounded. The answer separates into three parts that do not move
together.

\subsection{Shape: it does move, and this was expected to be the settled part}

The working expectation before the measurement was that the angular shape
would be settled: fixed by how the collapsed three-point cumulant decays
with transverse separation, and insensitive to the cutoff and to the
bispectrum model. That is not what the measurement shows, and it is worth
being explicit because it was the part assumed safe.

Normalising each curve to its own value at $0.5'$ and comparing over
$0.5'$ to $114'$:

\begin{center}
\begin{tabular}{lr}
\toprule
comparison & largest shape change \\
\midrule
BiHalofit against tree, both at $\elmax=960$ & $32\%$ \\
tree at $\elmax=7680$ against $\elmax=960$   & $62\%$ \\
tree at $\elmax=15360$ against $\elmax=960$  & $75\%$ \\
\bottomrule
\end{tabular}
\end{center}

The shape narrows as the cutoff rises. At $\gamma=13.6'$ the normalised
curve reads $0.44$ at $\elmax=960$ and $0.17$ at $\elmax=15360$. That
direction is physically sensible rather than alarming: a higher cutoff
admits shorter-wavelength modes, which decorrelate over shorter
separations, so the angular correlation narrows. The nonlinear model
narrows it slightly further at fixed cutoff, for the same reason.

What survives is the qualitative shape statement rather than a quantitative
one: the FK term decays monotonically with separation across the plotted
range, at every cutoff and for both bispectrum models, and it does so
without the flat plateau the deployed sampling produced. The rate of that
decay is not settled and should not be quoted as if it were.

\subsection{Amplitude: it does not converge to anything}

\begin{table}[h]
\centering
\begin{tabular}{rrrrrr}
\toprule
$\elmax$ & tree & BiHalofit & ratio & tree/O0 & hard $k$, near shell \\
\midrule
960 & $4.080\times10^{-6}$ & $6.603\times10^{-6}$ & 1.619 & 0.48\% & 7\,$h$/Mpc \\
1920 & $8.160\times10^{-6}$ & $21.630\times10^{-6}$ & 2.651 & 0.97\% & 13\,$h$/Mpc \\
3840 & $12.928\times10^{-6}$ & $67.662\times10^{-6}$ & 5.234 & 1.53\% & 26\,$h$/Mpc \\
7680 & $17.080\times10^{-6}$ & $181.821\times10^{-6}$ & 10.645 & 2.02\% & 52\,$h$/Mpc \\
15360 & $19.500\times10^{-6}$ & $347.230\times10^{-6}$ & 17.806 & 2.31\% & 105\,$h$/Mpc \\
\bottomrule
\end{tabular}
\caption{The FK contribution to $\xi_\kappa$ at $\gamma=0.5'$, tree level against BiHalofit, on the corrected grid, as a function of the multipole cutoff. Order-0 there is $8.443\times10^{-4}$. The last column is the largest wavenumber the bispectrum is asked for at the innermost shell; BiHalofit is calibrated to $k\lesssim10\,h$/Mpc. Produced by \code{code/rebuild/table\_nonlinear.py}.}
\label{tab:nonlinear}
\end{table}


This is sharper than expected, and it is the result that decides the
question. At the deployed cutoff the nonlinear bispectrum multiplies the
arcminute FK amplitude by $1.6$. That factor then grows with the cutoff
without settling: $2.7$, $5.2$, $10.6$, $17.8$ at
$\elmax=1920$, $3840$, $7680$, $15360$.

The tree-level and nonlinear cases behave qualitatively differently
(Figure~\ref{fig:cutoff}, middle panel). The tree curve flattens: its
successive increments are $2.00$, $1.58$, $1.32$, $1.14$, so it is
converging, toward roughly $2.5\%$ of Order-0. The BiHalofit curve does
not: its increments are $3.28$, $3.13$, $2.69$, $1.91$, and by
$\elmax=15360$ the FK term has reached $41\%$ of the Order-0 signal.

A term that is $41\%$ of the signal it corrects is not a perturbative
correction, and that number should not be read as a prediction. The reason
is structural rather than a matter of insufficient effort, and
Section~\ref{sec:uv} sets it out: the quantity being computed is a
coincident-point moment of a point-like beam, which has no finite value for
a nonlinear input. The cutoff is acting as an unstated regulator, so the
$41\%$ is a property of where the sum was stopped.

\subsection{Trustworthiness: the honest answer is negative}

Two separate limitations, and neither is removed by a better emulator.

\paragraph{BiHalofit is not independent of perturbation theory.} Its
three-halo term is built on the tree-level $F_2$ kernel with a fitted
effective spectrum. Only the one-halo term carries genuinely new,
simulation-calibrated information. So swapping tree for BiHalofit measures
how much the FK term grows once small-scale nonlinear power is included. It
does \emph{not} provide an independent check that the tree-level
calculation was right. The same holds for a bacco-style emulator: it shares
the gravity backbone and the whole projection chain.

\paragraph{The growth is concentrated at the lensing-kernel peak.} It is
worth asking which shells drive it, because the answer is not the obvious
one. Between $\elmax=960$ and $15360$ the tree vertex grows most at the
outer shells, by a factor $15$ at the source shell and only $1.5$ at the
innermost one: at the innermost shell the linear spectrum is already steep
enough that the integral had converged by $\elmax=960$. The nonlinear
vertex behaves differently, growing by a factor $69$ at
$\chi=3414$ Mpc, against $19$ at the source shell and $11$ at the
innermost. That intermediate shell is near the peak of the lensing
efficiency, and the wavenumber it reaches at the top cutoff is
$13\,h/$Mpc, which is exactly where the one-halo term flattens the spectrum
most. So the nonlinear growth is not confined to a corner of the integral
that the lensing kernel suppresses; it sits where the kernel is largest,
which is why the folded amplitude grows faster than the vertex at any
single shell.

\paragraph{And it happens where the model is not calibrated.} The last
column of the table is the largest wavenumber the bispectrum is asked for
at the innermost shell. At $\elmax=960$ it is $6.6\,h/$Mpc, inside
BiHalofit's calibration, and the factor $1.6$ measured there is a
defensible bound on the bispectrum-model contribution to the FK error
budget. By $\elmax=1920$ it is $13\,h/$Mpc, already outside. By $15360$ it
is $105\,h/$Mpc, which is inside single haloes, an order of magnitude
beyond the calibration and far into the regime where baryonic feedback
rather than gravity sets the matter spectrum. So the factor is trustworthy
exactly where it is smallest, and stops being trustworthy as it grows.

\subsection{The answer, in one sentence}

A nonlinear bispectrum gives a more realistic \emph{input} and an answer
that is not defined: the tree-level FK term converges to about $2.5\%$ of
Order-0 because the linear spectrum is steep enough to make the
coincident-point moment finite, while the nonlinear one has no limit
without a smoothing scale the formalism does not supply
(Section~\ref{sec:uv}), reaching $41\%$ of Order-0 at a cutoff where its
dominant modes sit at $100\,h/$Mpc. The defensible presentation is the
amplitude as a function of the small-scale cutoff, with the model
comparison quoted only where the cutoff keeps the dominant modes inside the
calibrated range, which is the deployed cutoff and not much beyond it.

This is a negative answer to the question the programme started from, and
it is worth stating as such rather than as a caveat attached to a quoted
nonlinear number.

\subsection{What would actually validate the model}

Neither BiHalofit nor an emulator can, for the reason above. Two things
could, and neither has been done:

\begin{enumerate}
\item A convergence-bispectrum comparison: predict
$B_\kappa(\ell_1,\ell_2,\ell_3)$ with the same machinery and compare
against direct measurements on public ray-traced lightcones. This tests the
Poisson factor, the screen-Hessian normalisation, the Limber projection and
the measure together. It cannot reach the coincident limit the FK term
samples, so it does not close the gap, but nothing available gets closer.
\item A direct measurement of the collapsed three-point cumulant of the
driving fields on N-body lightcones, which does close it.
\end{enumerate}
